\documentclass[11pt]{article}
\usepackage[margin=1in]{geometry}
\usepackage{amsmath, amssymb, amsthm}
\usepackage{enumitem}
\usepackage{booktabs}
\usepackage{microtype}
\usepackage{tikz}
\usetikzlibrary{calc}
\usepackage{hyperref}

\newcommand{\X}[2]{X_{#1#2}}
\newcommand{\Z}{\mathcal{Z}}
\newtheorem{theorem}{Theorem}
\newtheorem{lemma}{Lemma}
\newtheorem{proposition}{Proposition}
\theoremstyle{definition}
\newtheorem{definition}{Definition}
\newtheorem*{remark}{Remark}

\title{Locality emerges from hidden zeros in \texorpdfstring{$\mathrm{tr}(\phi^3)$}{tr(phi^3)}:\\
a by-hand proof at four, five, and six points}
\author{Jony}
\date{\today}

\begin{document}
\maketitle

\begin{abstract}
We prove that the cyclic hidden one-zeros of color-ordered $\mathrm{tr}(\phi^3)$
tree amplitudes force every non-local coefficient of the general rational
ansatz to vanish, at $n=4, 5, 6$ external legs. The proof is entirely by hand:
we use only the kinematic mesh identity, geometric series Laurent expansions,
and the fact that rational functions in independent variables vanish iff
each coefficient on a basis of monomials vanishes.
The mechanism is geometric: each cyclic zero kills exactly the non-local
terms whose chord factors involve the ``enclosed vertex'' of a designated
short diagonal (the \emph{special variable}). Every non-local coefficient is
killed by at least one zero --- and as it happens, by exactly two.
\end{abstract}

\tableofcontents

%================================================================
\section{Setup}

\subsection{Chords, the ansatz, and locality}

Label the external legs cyclically by $1,\ldots,n$. The planar Mandelstam
variables are
\[
\X{i}{j} \;=\; (p_i+p_{i+1}+\cdots+p_{j-1})^2,\qquad 1\le i<j\le n,\ j-i\ge 2,\ (i,j)\ne(1,n).
\]
Geometrically, $\X{i}{j}$ is the chord from vertex $i$ to vertex $j$ of the
convex $n$-gon. There are $n(n-3)/2$ such chords ($2,5,9$ for $n=4,5,6$).

A \emph{triangulation} of the $n$-gon is a maximal set of non-crossing chords;
it has $n-3$ chords. The tree amplitude is the sum over triangulations of the
product of chord propagators,
\[
A_n \;=\;\sum_{T\text{ triangulation}}\ \prod_{(i,j)\in T}\frac{1}{\X{i}{j}}.
\]
The number of triangulations is the Catalan number $C_{n-2}$ ($2, 5, 14$
for $n=4,5,6$).

We consider the most general rational function $B$ with the correct dimension
($-2(n-3)$) and only planar simple poles:
\[
B \;=\; \sum_{\text{multisets }(a_1,\ldots,a_{n-3})}
\frac{a_{a_1\ldots a_{n-3}}}{X_{a_1}\cdots X_{a_{n-3}}}.
\]
The sum is over multisets (unordered with repetition) of size $n-3$ from
the set of chord indices. There are $\binom{d+n-4}{n-3}$ such multisets,
where $d=n(n-3)/2$ is the number of chords. We get $2, 15, 165$ multisets
for $n=4,5,6$.

A multiset is \emph{local} if its chords form a triangulation; the remaining
multisets are \emph{non-local} (either repeated chords or two chords that
cross). Locality is the statement that all non-local coefficients vanish.

\subsection{The kinematic mesh identity}

The non-planar two-particle invariants $c_{a,b}=2p_a\cdot p_b$ ($a\ne b$,
non-adjacent) are not independent of the planar variables. They satisfy
\[
\boxed{\ c_{a,b} \;=\; \X{a}{b+1} + \X{a+1}{b} - \X{a}{b} - \X{a+1}{b+1}\ }
\]
(adjacent chords like $\X{a}{a+1}$ are zero by masslessness $p_a^2=0$).
Setting $c_{a,b}=0$ thus imposes a linear relation on the planar chords.

\subsection{The hidden one-zeros}

For each leg $r\in\{1,\ldots,n\}$, the \emph{hidden one-zero} $\Z_r$ is
defined by setting all $c$-invariants involving leg $r$ to zero:
\[
\Z_r:\qquad c_{r,r+2}=c_{r,r+3}=\cdots=c_{r,r+n-2}=0\ \pmod n.
\]
There are $n-3$ such constraints. The starting fact (proved elsewhere; here
taken as given) is that the tree amplitude $A_n$ vanishes on each $\Z_r$.
Our problem is the \emph{converse} for the general ansatz: assuming
$B|_{\Z_r}=0$ for all $r$, show that all non-local $a$-coefficients vanish.

\subsection{The geometric heart: enclosed vertices}

Of the two arcs into which a chord $\X{i}{j}$ divides the polygon's
boundary, the \emph{enclosed arc} is the shorter one. The
\emph{enclosed vertex set} is
\[
E(\X{i}{j}) \;=\; \{\text{vertices strictly between $i$ and $j$ on the shorter arc}\}.
\]
We call $|E(\X{i}{j})|$ the \emph{class} of the chord:
class 1 chords (``short diagonals'') enclose $1$ vertex, class 2 (``long
diagonals'') enclose $2$, etc.

\begin{lemma}[Crossing $=$ enclosed-vertex sweep]
\label{lem:crossing}
Two chords $\X{i}{j}$ and $\X{r}{s}$ of a convex polygon cross in the
interior iff exactly one endpoint of $\X{r}{s}$ lies in $E(\X{i}{j})$
(equivalently in the enclosed arc of $\X{i}{j}$), the other in the outer arc.
\end{lemma}

\begin{proof}
Two chords of a convex polygon cross iff their four endpoints alternate around
the boundary. Fix $\X{i}{j}$ and consider the two arcs it cuts off (one
``enclosed,'' one ``outer''). The endpoints of $\X{r}{s}$ either both lie in
the same arc (no alternation, no crossing) or one lies in each (alternation,
crossing).
\end{proof}

\paragraph{Consequence.}
Every non-local pair of chords --- i.e.\ a crossing pair --- has the form
``one chord with enclosed vertex $v$, another chord starting at $v$.'' The
hidden zeros, as we'll see, are tailor-built to kill these.

%================================================================
\section{Four points: trivial warm-up}

The square has just two chords, $\X{1}{3}$ and $\X{2}{4}$. The ansatz
\[
B = \frac{a_{\X{1}{3}}}{\X{1}{3}} + \frac{a_{\X{2}{4}}}{\X{2}{4}}
\]
has $2$ coefficients, both \emph{local} (each chord by itself is a triangulation).
There are no non-local terms to kill. A single zero $\Z_1$ gives
$c_{1,3}=\X{1}{4}+\X{2}{3}-\X{1}{3}-\X{2}{4}=-\X{1}{3}-\X{2}{4}=0$, hence
$\X{2}{4}=-\X{1}{3}$, and $B|_{\Z_1}=(a_{\X{1}{3}}-a_{\X{2}{4}})/\X{1}{3}=0$
forces $a_{\X{1}{3}}=a_{\X{2}{4}}$. This is unitarity (equal residues on both
channels), not locality. We move on to $n=5$, where the question becomes
nontrivial.

%================================================================
\section{Five points: complete proof}
\label{sec:n5}

\subsection{The pentagon's chords and crossings}

The five chords of the pentagon, listed with their enclosed-vertex sets:
\begin{center}
\begin{tabular}{cccccc}
\toprule
Chord & $\X{1}{3}$ & $\X{1}{4}$ & $\X{2}{4}$ & $\X{2}{5}$ & $\X{3}{5}$ \\
$E(\cdot)$ & $\{2\}$ & $\{5\}$ & $\{3\}$ & $\{1\}$ & $\{4\}$ \\
\bottomrule
\end{tabular}
\end{center}

Every chord has class $1$ (encloses one vertex). At $n=5$ there are no
class-2 chords yet; they appear starting at $n=6$.

The five \emph{crossing pairs} (verified by Lemma \ref{lem:crossing}): for each
pair, one chord's enclosed vertex is an endpoint of the other.
\begin{center}
\begin{tabular}{cccc}
\toprule
Pair & $E$-witnesses & Pair & $E$-witnesses \\
\midrule
$\X{1}{3}\times\X{2}{4}$ & $2\in\X{2}{4},\ 3\in\X{1}{3}^c$ &
$\X{1}{4}\times\X{3}{5}$ & $5\in\X{3}{5}^c,\ 4\in\X{1}{4}^c$ \\
$\X{1}{3}\times\X{2}{5}$ & $2\in\X{2}{5},\ 1\in\X{1}{3}^c$ &
$\X{2}{4}\times\X{3}{5}$ & $3\in\X{3}{5}^c,\ 4\in\X{2}{4}^c$ \\
$\X{1}{4}\times\X{2}{5}$ & $5\in\X{2}{5}^c,\ 1\in\X{1}{4}^c$ & & \\
\bottomrule
\end{tabular}
\end{center}

The other $5$ unordered pairs of distinct chords share an endpoint (so
non-crossing) and are exactly the $5$ fan triangulations.

\subsection{The 15-coefficient ansatz}

With $n-3=2$ chord factors per monomial, the ansatz is
\[
B = \sum_{1\le i\le j\le 5}\frac{a_{ij}}{X_i X_j},\quad\text{15 coefficients},
\]
where the indexing runs over multisets of chord-indices. Of the $15$ coefficients:
\begin{itemize}[nosep,leftmargin=2em]
\item $5$ \emph{tree triples}: one for each fan triangulation.
\item $5$ \emph{double poles}: $a_{\X{i}{j}\X{i}{j}}$, one per chord.
\item $5$ \emph{crossing pairs}: one $a_{\X{i}{j}\X{k}{\ell}}$ per crossing.
\end{itemize}
The latter $10$ are the non-local coefficients we must kill.

\subsection{The cyclic hidden one-zeros, derived by hand}

We derive each $\Z_r$ from $c_{r,r+2}=c_{r,r+3}=0$ using the mesh identity.

\paragraph{$\Z_1$:} $c_{1,3}=c_{1,4}=0$.
\begin{align*}
c_{1,3}&=\X{1}{4}+\X{2}{3}-\X{1}{3}-\X{2}{4}=\X{1}{4}-\X{1}{3}-\X{2}{4}=0,\\
c_{1,4}&=\X{1}{5}+\X{2}{4}-\X{1}{4}-\X{2}{5}=\X{2}{4}-\X{1}{4}-\X{2}{5}=0.
\end{align*}
(We used $\X{2}{3}=\X{1}{5}=0$, since $\X{a}{a+1}$ corresponds to the
massless particle squared.) Solving for $\X{2}{4}$ and $\X{2}{5}$ (the two
chords from the enclosed vertex $2$ of $\X{1}{3}$):
\begin{equation}
\Z_1:\qquad\boxed{\X{2}{4}=\X{1}{4}-\X{1}{3},\qquad \X{2}{5}=-\X{1}{3}.}
\end{equation}
The chord $\X{1}{3}$ appears on the right of \emph{both} substitutions: it
is the \textbf{special variable} of $\Z_1$.

\paragraph{$\Z_2$:} $c_{2,4}=c_{2,5}=0$.
$c_{2,4}=\X{2}{5}+\X{3}{4}-\X{2}{4}-\X{3}{5}=\X{2}{5}-\X{2}{4}-\X{3}{5}=0$.
$c_{2,5}=\X{2}{6}+\X{3}{5}-\X{2}{5}-\X{3}{6}$, where $\X{2}{6}=\X{2}{1}=0$
and $\X{3}{6}=\X{3}{1}=\X{1}{3}$, so $c_{2,5}=\X{3}{5}-\X{2}{5}-\X{1}{3}=0$.
Solving:
\begin{equation}
\Z_2:\qquad\boxed{\X{3}{5}=\X{2}{5}-\X{2}{4},\qquad \X{1}{3}=-\X{2}{4}}, \quad\text{special }=\X{2}{4}.
\end{equation}

\paragraph{$\Z_3$:} $c_{3,5}=c_{3,1}=0$ (cyclic).
$c_{3,5}=\X{3}{6}+\X{4}{5}-\X{3}{5}-\X{4}{6}=\X{1}{3}-\X{3}{5}-\X{1}{4}=0$, so
$\X{1}{4}=\X{1}{3}-\X{3}{5}$.
$c_{3,1}=\X{3}{2}+\X{4}{1}-\X{3}{1}-\X{4}{2}=\X{1}{4}-\X{1}{3}-\X{2}{4}=0$, so
$\X{2}{4}=\X{1}{4}-\X{1}{3}=-\X{3}{5}$.
\begin{equation}
\Z_3:\qquad\boxed{\X{1}{4}=\X{1}{3}-\X{3}{5},\qquad \X{2}{4}=-\X{3}{5}}, \quad\text{special }=\X{3}{5}.
\end{equation}

\paragraph{$\Z_4$:} $c_{4,1}=c_{4,2}=0$.
$c_{4,1}=\X{4}{2}+\X{5}{1}-\X{4}{1}-\X{5}{2}=\X{2}{4}-\X{1}{4}-\X{2}{5}=0$, so $\X{2}{5}=\X{2}{4}-\X{1}{4}$.
$c_{4,2}=\X{4}{3}+\X{5}{2}-\X{4}{2}-\X{5}{3}=\X{2}{5}-\X{2}{4}-\X{3}{5}=0$, so $\X{3}{5}=\X{2}{5}-\X{2}{4}=-\X{1}{4}$.
\begin{equation}
\Z_4:\qquad\boxed{\X{2}{5}=\X{2}{4}-\X{1}{4},\qquad \X{3}{5}=-\X{1}{4}}, \quad\text{special }=\X{1}{4}.
\end{equation}

\paragraph{$\Z_5$:} $c_{5,2}=c_{5,3}=0$.
$c_{5,2}=\X{5}{3}+\X{1}{2}-\X{5}{2}-\X{1}{3}=\X{3}{5}-\X{2}{5}-\X{1}{3}=0$, so $\X{1}{3}=\X{3}{5}-\X{2}{5}$.
$c_{5,3}=\X{5}{4}+\X{1}{3}-\X{5}{3}-\X{1}{4}=\X{1}{3}-\X{3}{5}-\X{1}{4}=0$, so $\X{1}{4}=\X{1}{3}-\X{3}{5}=-\X{2}{5}$.
\begin{equation}
\Z_5:\qquad\boxed{\X{1}{3}=\X{3}{5}-\X{2}{5},\qquad \X{1}{4}=-\X{2}{5}}, \quad\text{special }=\X{2}{5}.
\end{equation}

\subsection{The pattern: every zero is a rotation of the previous}

The structure of every zero is identical:
\begin{itemize}[nosep,leftmargin=2em]
\item \emph{One} special chord $X_*$ appears in both substitutions.
\item \emph{Two} substituted chords (the chords from the enclosed vertex of $X_*$).
  One has the form ``$X_s = X_f - X_*$'' with a non-zero \textbf{companion} $X_f$;
  the other is ``bare'': $X_s = -X_*$ (companion zero).
\item \emph{Two} free chords $F$ remain. One of them is the
  companion of the non-bare substitution.
\end{itemize}

\subsection{The kill technique, in full detail at \texorpdfstring{$\Z_1$}{Z1}}

We work out $\Z_1$ in full and verify the kills by hand. The other four
zeros follow by cyclic rotation; their kills can be read off.

For $\Z_1$, the special is $X_*=\X{1}{3}$, the substitutions are
$\X{2}{4}=\X{1}{4}-\X{1}{3}$ (non-bare, companion $\X{1}{4}$) and
$\X{2}{5}=-\X{1}{3}$ (bare). The free set is $F_1=\{\X{1}{4},\X{3}{5}\}$.

After substitution, $B|_{\Z_1}$ is a rational function of three variables:
$\X{1}{3},\X{1}{4},\X{3}{5}$. The condition $B|_{\Z_1}=0$ means this
rational function vanishes identically. We extract single-coefficient
kills by Laurent-expanding in the special variable $\X{1}{3}$, treating
$\X{1}{4},\X{3}{5}$ as fixed.

\paragraph{Step 1: Order $\X{1}{3}^0$ (the limit $\X{1}{3}\to\infty$).}

As $\X{1}{3}\to\infty$ with $\X{1}{4},\X{3}{5}$ held fixed:
\[
\X{2}{4}=\X{1}{4}-\X{1}{3}\to-\infty,\quad \X{2}{5}=-\X{1}{3}\to-\infty,
\]
so $1/\X{2}{4}\to 0$ and $1/\X{2}{5}\to 0$. Thus every monomial of $B$
containing any factor of $\X{1}{3}$, $\X{2}{4}$, or $\X{2}{5}$ vanishes in
the limit. The surviving monomials use only the two free chords $\X{1}{4},\X{3}{5}$:
\[
\lim_{\X{1}{3}\to\infty} B|_{\Z_1}
\;=\;
\frac{a_{\X{1}{4}\X{1}{4}}}{\X{1}{4}^2} +
\frac{a_{\X{3}{5}\X{3}{5}}}{\X{3}{5}^2} +
\frac{a_{\X{1}{4}\X{3}{5}}}{\X{1}{4}\,\X{3}{5}}.
\]
This must vanish identically as a function of two independent variables
$\X{1}{4},\X{3}{5}$. The three rational functions $1/\X{1}{4}^2$, $1/\X{3}{5}^2$,
$1/(\X{1}{4}\X{3}{5})$ are linearly independent. Hence each coefficient must be zero:
\[
\boxed{\quad a_{\X{1}{4}\X{1}{4}}=a_{\X{3}{5}\X{3}{5}}=a_{\X{1}{4}\X{3}{5}}=0.\quad}
\quad\text{($\Z_1$, three order-0 kills)}
\]

\paragraph{Step 2: Order $\X{1}{3}^{-1}$ (coefficient of $1/\X{1}{3}$).}

Expand the substituted chords as geometric series:
\[
\frac{1}{\X{2}{4}}
=\frac{1}{\X{1}{4}-\X{1}{3}}
=-\frac{1}{\X{1}{3}}\cdot\frac{1}{1-\X{1}{4}/\X{1}{3}}
=-\frac{1}{\X{1}{3}}-\frac{\X{1}{4}}{\X{1}{3}^2}-\frac{\X{1}{4}^2}{\X{1}{3}^3}-\cdots,
\]
\[
\frac{1}{\X{2}{5}}=-\frac{1}{\X{1}{3}}\quad\text{(no tail).}
\]
Tracking the $1/\X{1}{3}$ coefficient of each monomial, we find
\[
\bigl(B|_{\Z_1}\bigr)_{1/\X{1}{3}}
=
\frac{a_{\X{1}{3}\X{1}{4}}-a_{\X{2}{4}\X{1}{4}}-a_{\X{2}{5}\X{1}{4}}}{\X{1}{4}}
+\frac{a_{\X{1}{3}\X{3}{5}}-a_{\X{2}{4}\X{3}{5}}-a_{\X{2}{5}\X{3}{5}}}{\X{3}{5}}.
\]
Setting this to zero and using linear independence of $1/\X{1}{4}$ and
$1/\X{3}{5}$:
\begin{align*}
a_{\X{1}{3}\X{1}{4}}&=a_{\X{2}{4}\X{1}{4}}+a_{\X{2}{5}\X{1}{4}},\\
a_{\X{1}{3}\X{3}{5}}&=a_{\X{2}{4}\X{3}{5}}+a_{\X{2}{5}\X{3}{5}}.
\end{align*}
These are \emph{multi-term} relations, not single-coefficient kills.

\paragraph{Step 3: Order $\X{1}{3}^{-2}$ (coefficient of $1/\X{1}{3}^2$).}

Tabulating contributions monomial-by-monomial. The crucial new piece is
the monomial $a_{\X{2}{4}\X{3}{5}}/(\X{2}{4}\X{3}{5})$: its $1/\X{1}{3}^2$
contribution is
\[
-\frac{a_{\X{2}{4}\X{3}{5}}\,\X{1}{4}}{\X{1}{3}^{\,2}\,\X{3}{5}},
\]
producing a free-variable monomial $\X{1}{4}/\X{3}{5}$. \emph{No other
monomial of $B$ contributes a $\X{1}{4}/\X{3}{5}$ structure at order
$\X{1}{3}^{-2}$}: all other monomials at this order produce either
constants (from the leading $-1/\X{1}{3}$ piece of the substituted chord
in monomials like $1/(\X{1}{3}\X{2}{4})$, $1/\X{2}{4}^2$, etc.) or other
free-variable monomials.

Setting the coefficient of $\X{1}{4}/\X{3}{5}$ to zero in
$(B|_{\Z_1})_{1/\X{1}{3}^2}=0$:
\[
\boxed{a_{\X{2}{4}\X{3}{5}}=0\qquad\text{($\Z_1$, single order-2 kill).}}
\]

(The constant part of $(B|_{\Z_1})_{1/\X{1}{3}^2}$ gives a multi-term
equation involving $7$ coefficients, which we don't pursue here.)

\paragraph{Step 4: orders $\X{1}{3}^{-3}, \X{1}{3}^{-4}, \ldots$.}

A direct check shows that no new single-coefficient kills arise. The same
monomial $\X{1}{4}/\X{3}{5}$ recurs (giving $a_{\X{2}{4}\X{3}{5}}=0$
redundantly), and new monomials produce multi-term equations.

\paragraph{Summary of $\Z_1$:}
\[
\text{Single-coefficient kills from }\Z_1:\quad
\boxed{a_{\X{1}{4}\X{1}{4}},\ a_{\X{3}{5}\X{3}{5}},\ a_{\X{1}{4}\X{3}{5}}}\ \text{(order 0)},\quad
\boxed{a_{\X{2}{4}\X{3}{5}}}\ \text{(order 2)}.
\]
Four kills total.

\subsection{Cycling: kills from \texorpdfstring{$\Z_2,\ldots,\Z_5$}{Z2..Z5}}

By cyclic symmetry, each $\Z_r$ has the same structure as $\Z_1$, rotated.
Reading off:

\begin{center}
\renewcommand{\arraystretch}{1.2}
\begin{tabular}{c|c|c}
\toprule
Zero & Order-0 kills (multisets in $F_r$) & Order-2 kill \\
\midrule
$\Z_1$ ($X_*=\X{1}{3}$) & $a_{\X{1}{4}\X{1}{4}},\ a_{\X{3}{5}\X{3}{5}},\ a_{\X{1}{4}\X{3}{5}}$ & $a_{\X{2}{4}\X{3}{5}}$\\
$\Z_2$ ($X_*=\X{2}{4}$) & $a_{\X{1}{4}\X{1}{4}},\ a_{\X{2}{5}\X{2}{5}},\ a_{\X{1}{4}\X{2}{5}}$ & $a_{\X{1}{4}\X{3}{5}}$\\
$\Z_3$ ($X_*=\X{3}{5}$) & $a_{\X{1}{3}\X{1}{3}},\ a_{\X{2}{5}\X{2}{5}},\ a_{\X{1}{3}\X{2}{5}}$ & $a_{\X{1}{4}\X{2}{5}}$\\
$\Z_4$ ($X_*=\X{1}{4}$) & $a_{\X{1}{3}\X{1}{3}},\ a_{\X{2}{4}\X{2}{4}},\ a_{\X{1}{3}\X{2}{4}}$ & $a_{\X{1}{3}\X{2}{5}}$\\
$\Z_5$ ($X_*=\X{2}{5}$) & $a_{\X{2}{4}\X{2}{4}},\ a_{\X{3}{5}\X{3}{5}},\ a_{\X{2}{4}\X{3}{5}}$ & $a_{\X{1}{3}\X{2}{4}}$\\
\bottomrule
\end{tabular}
\end{center}

\subsection{Union: every non-local coefficient is killed}

\paragraph{Double poles (5 total):} each is killed at order 0 by exactly two zones.
\begin{itemize}[nosep,leftmargin=2em]
\item $a_{\X{1}{3}^2}$: by $\Z_3$ and $\Z_4$.
\item $a_{\X{1}{4}^2}$: by $\Z_1$ and $\Z_2$.
\item $a_{\X{2}{4}^2}$: by $\Z_4$ and $\Z_5$.
\item $a_{\X{2}{5}^2}$: by $\Z_2$ and $\Z_3$.
\item $a_{\X{3}{5}^2}$: by $\Z_1$ and $\Z_5$.
\end{itemize}

\paragraph{Crossing pairs (5 total):} each is killed once at order 0 and once at order 2.
\begin{itemize}[nosep,leftmargin=2em]
\item $a_{\X{1}{3}\X{2}{4}}$: order 0 by $\Z_4$, order 2 by $\Z_5$.
\item $a_{\X{1}{3}\X{2}{5}}$: order 0 by $\Z_3$, order 2 by $\Z_4$.
\item $a_{\X{1}{4}\X{2}{5}}$: order 0 by $\Z_2$, order 2 by $\Z_3$.
\item $a_{\X{1}{4}\X{3}{5}}$: order 0 by $\Z_1$, order 2 by $\Z_2$.
\item $a_{\X{2}{4}\X{3}{5}}$: order 0 by $\Z_5$, order 2 by $\Z_1$.
\end{itemize}

\paragraph{Tree triples (5 total) are never killed.} For each fan triangulation
$a_{\X{i}{j}\X{j}{k}}$ (chords sharing endpoint $j$), at every zone $\Z_r$
at least one of the two indices is either the special $X_*$, the bare
substituted chord, or one of the substituted chords adjacent to the special
in a way that breaks the conditions for both order-0 and order-2 kills.
Equivalently: the tree amplitude $A_5$ vanishes on each $\Z_r$, so every
Laurent coefficient of $A_5|_{\Z_r}$ vanishes; the tree triples
in $A_5$ form the kernel of all our conditions.

\begin{theorem}[Locality at five points]
The five hidden one-zeros of the pentagon force every non-local coefficient
of the general rational ansatz $B$ to vanish. The surviving 5 coefficients
are precisely the tree triangulations.
\end{theorem}

\begin{proof}
By the explicit calculation at $\Z_1$ above (Steps 1--4) and cyclic rotation
to obtain $\Z_2,\ldots,\Z_5$, the union of single-coefficient kills covers
all 10 non-local coefficients (5 double poles, 5 crossing pairs), as
tabulated. The 5 tree triples are never killed. \qed
\end{proof}

\subsection{The geometric pattern revealed}
\label{sec:geometric-pattern-n5}

Look at the table of which zone kills each crossing pair at which order.
A clean pattern emerges: every crossing pair is killed by exactly two zones,
and these zones are \emph{adjacent} in the cyclic ordering. The order-0
kill comes from the zone where both chords lie in the free set; the order-2
kill comes from the next zone over, where the special variable becomes one
of the chord-pair's neighbors via the enclosed-vertex correspondence.

This is the \textbf{enclosed-vertex correspondence in action}: the hidden
zeros are tightly linked to the crossings of the polygon, and every
crossing has a precise pair of killing zones determined by the
enclosed-vertex structure.

%================================================================
\section{Six points: the hexagon case}
\label{sec:n6}

The same technique works at $n=6$, with two new features:
\begin{itemize}[nosep,leftmargin=2em]
\item The hexagon has \emph{two classes} of chords (short and long diagonals).
\item Each zero now has \emph{three} substitutions, of which \emph{two} are
non-bare. This produces an additional ``order 4'' layer of kills coming
from monomials with two non-bare substituted indices.
\end{itemize}

\subsection{Setup at $n=6$}

The hexagon's 9 chords:
\begin{center}
\begin{tabular}{ll}
\textbf{Short diagonals} (class 1, enclose 1 vertex): &
$\X{1}{3},\X{2}{4},\X{3}{5},\X{4}{6},\X{1}{5},\X{2}{6}$ \\
\textbf{Long diagonals} (class 2, enclose 2 vertices): &
$\X{1}{4},\X{2}{5},\X{3}{6}$
\end{tabular}
\end{center}

Cyclic shorthand:
\[
x_1=\X{1}{3},\ x_2=\X{2}{4},\ x_3=\X{3}{5},\ x_4=\X{4}{6},\ x_5=\X{1}{5},\ x_6=\X{2}{6};
\]
\[
y_1=\X{1}{4},\ y_2=\X{2}{5},\ y_3=\X{3}{6}.
\]

Enclosed-vertex sets:
\[
E(x_r) = \{r+1\}\bmod 6,\qquad E(y_r) = \{r+1, r+2\}\bmod 6.
\]

\paragraph{Ansatz.} Each monomial has $n-3=3$ chord factors. The number of
multisets of size 3 from 9 chords is $\binom{9+2}{3}=165$. Of these, $14$
are tree triples ($C_4=14$ triangulations of the hexagon), $151$ are
non-local.

\subsection{The six zeros, derived by hand}

Each $\Z_r$ is defined by $c_{r,r+2}=c_{r,r+3}=c_{r,r+4}=0$ (three constraints).
For $\Z_1$:

$c_{1,3}=\X{1}{4}+\X{2}{3}-\X{1}{3}-\X{2}{4}=y_1-x_1-x_2=0\Rightarrow x_2=y_1-x_1$.

$c_{1,4}=\X{1}{5}+\X{2}{4}-\X{1}{4}-\X{2}{5}=x_5+x_2-y_1-y_2=0$.
Using $x_2=y_1-x_1$: $x_5+y_1-x_1-y_1-y_2=x_5-x_1-y_2=0$, so $y_2=x_5-x_1$.

$c_{1,5}=\X{1}{6}+\X{2}{5}-\X{1}{5}-\X{2}{6}=0+y_2-x_5-x_6=0$.
Using $y_2=x_5-x_1$: $-x_1-x_6=0$, so $x_6=-x_1$.

So $\Z_1$: $x_2=y_1-x_1$, $y_2=x_5-x_1$, $x_6=-x_1$, with special $x_*=x_1=\X{1}{3}$.
The substituted chords $\{x_2,y_2,x_6\}=\{\X{2}{4},\X{2}{5},\X{2}{6}\}$ are
exactly the chords from the enclosed vertex $2$ of $\X{1}{3}$.

By identical derivations (cyclic), the six zeros are:

\begin{center}
\renewcommand{\arraystretch}{1.2}
\begin{tabular}{c|c|c|c}
\toprule
Zero & Special & Substitutions & Free $F_r$ (excluding spec.)\\
\midrule
$\Z_1$ & $x_1$ & $x_2=y_1-x_1$, $y_2=x_5-x_1$, $x_6=-x_1$ & $\{x_3,x_4,x_5,y_1,y_3\}$\\
$\Z_2$ & $x_2$ & $x_3=y_2-x_2$, $y_3=x_6-x_2$, $x_1=-x_2$ & $\{x_4,x_5,x_6,y_1,y_2\}$\\
$\Z_3$ & $x_3$ & $x_4=y_3-x_3$, $y_1=x_1-x_3$, $x_2=-x_3$ & $\{x_1,x_5,x_6,y_2,y_3\}$\\
$\Z_4$ & $x_4$ & $x_5=y_1-x_4$, $y_2=x_2-x_4$, $x_3=-x_4$ & $\{x_1,x_2,x_6,y_1,y_3\}$\\
$\Z_5$ & $x_5$ & $x_6=y_2-x_5$, $y_3=x_3-x_5$, $x_4=-x_5$ & $\{x_1,x_2,x_3,y_1,y_2\}$\\
$\Z_6$ & $x_6$ & $x_1=y_3-x_6$, $y_1=x_4-x_6$, $x_5=-x_6$ & $\{x_2,x_3,x_4,y_2,y_3\}$\\
\bottomrule
\end{tabular}
\end{center}

In every zone:
\begin{itemize}[nosep,leftmargin=2em]
\item One special $x_*$ (a short diagonal).
\item Three substituted chords from the enclosed vertex of $x_*$. Of these:
  \emph{two are non-bare} (one short, one long, with companions in $F_r$),
  and \emph{one is bare} ($=-x_*$).
\item Five free chords $F_r$ (three short, two long).
\end{itemize}

\subsection{The kill technique at \texorpdfstring{$\Z_1$}{Z1}: three layers}

For $\Z_1$: $x_*=x_1$, non-bare substitutions
\[
x_2=y_1-x_1\ \text{(companion $f_{x_2}=y_1$)},\quad y_2=x_5-x_1\ \text{(companion $f_{y_2}=x_5$)},
\]
bare: $x_6=-x_1$. Free: $F_1=\{x_3,x_4,x_5,y_1,y_3\}$ (5 chords).

The Laurent expansions:
\[
\frac{1}{x_2}=-\frac{1}{x_1}-\frac{y_1}{x_1^2}-\frac{y_1^2}{x_1^3}-\cdots,
\quad
\frac{1}{y_2}=-\frac{1}{x_1}-\frac{x_5}{x_1^2}-\frac{x_5^2}{x_1^3}-\cdots,
\quad
\frac{1}{x_6}=-\frac{1}{x_1}.
\]

\paragraph{Layer 1: Order $x_1^0$ ($x_1\to\infty$ limit). [\textbf{35 kills}]}

As $x_1\to\infty$, all of $x_2,y_2,x_6$ tend to $\pm\infty$, so $1/x_2,1/y_2,1/x_6\to 0$.
Surviving monomials of $B$ are those with all three indices in $F_1$:
\[
\bigl(B|_{\Z_1}\bigr)\bigr|_{x_1\to\infty} = \sum_{(a,b,c)\subseteq F_1}\frac{a_{abc}}{X_a X_b X_c}.
\]
The number of multisets $(a,b,c)\subseteq F_1$ (size 3 from 5 elements) is $\binom{5+2}{3}=35$.
Each $1/(X_aX_bX_c)$ is a distinct rational monomial in 5 free variables;
they are linearly independent. Hence each $a_{abc}=0$:
\[
\boxed{\text{35 single-term kills at order 0: all }a_{abc}\text{ with }\{X_a,X_b,X_c\}\subseteq F_1.}
\]

\paragraph{Layer 2: Order $x_1^{-2}$ (coefficient of $1/x_1^2$). [\textbf{20 new kills}]}

Consider monomials of $B$ with exactly one factor among the non-bare
substituted chords $\{x_2,y_2\}$, the other two factors in $F_1$. We argue
by the ``unique source'' principle.

\textbf{Sub-case A: monomial $a_{x_2 X_a X_b}/(x_2 X_a X_b)$ with $X_a,X_b\in F_1$.}
Substituting $1/x_2=-1/x_1-y_1/x_1^2-\cdots$, the $1/x_1^2$ piece is
$-y_1/(x_1^2 X_a X_b)$, contributing
\[
-\frac{a_{x_2 X_a X_b}\cdot y_1}{X_a X_b}\quad\text{to the $1/x_1^2$ coefficient.}
\]

\emph{When is the monomial $y_1/(X_a X_b)$ a unique source?} It must not arise
from any other term. The only other potential source would be from
$1/y_2$'s tail, but $1/y_2$'s order-2 tail term is $x_5/x_1^2$, not $y_1$;
so $y_1/(X_a X_b)$ uniquely identifies the $a_{x_2 X_a X_b}$ origin
\emph{provided} neither of $X_a, X_b$ is $y_1$ (otherwise the $y_1$ in
the numerator ``cancels'' with the $1/y_1$ in the denominator, producing
a degenerate monomial that mixes with other contributions).

Hence the unique-source condition: $\{X_a,X_b\}\subseteq F_1\setminus\{y_1\}=\{x_3,x_4,x_5,y_3\}$.
Number of such multisets of size 2 from 4 elements: $\binom{4+1}{2}=10$.

\[
\boxed{\text{10 single-term kills at order 2: all }a_{x_2 X_a X_b}=0\text{ with }\{X_a,X_b\}\subseteq\{x_3,x_4,x_5,y_3\}.}
\]

\textbf{Sub-case B: monomial $a_{y_2 X_a X_b}/(y_2 X_a X_b)$ with $X_a,X_b\in F_1$.}
Same argument with $y_2$'s companion $x_5$ instead of $y_1$. Unique-source
condition: $\{X_a,X_b\}\subseteq F_1\setminus\{x_5\}=\{x_3,x_4,y_1,y_3\}$.
Another $\binom{4+1}{2}=10$ kills.

\[
\boxed{\text{10 more single-term kills at order 2: all }a_{y_2 X_a X_b}=0\text{ with }\{X_a,X_b\}\subseteq\{x_3,x_4,y_1,y_3\}.}
\]

(Sub-case for $x_6$: bare substitution, contributes $0$ to the order-2
companion structure since $1/x_6$ has no tail.)

\paragraph{Layer 3: Order $x_1^{-4}$ (coefficient of $1/x_1^4$). [\textbf{3 new kills}]}

Consider monomials of $B$ with \emph{both} non-bare substitutions as
factors: $a_{x_2 y_2 X_c}/(x_2 y_2 X_c)$ with $X_c\in F_1$.

Use the product Laurent expansion:
\[
\frac{1}{x_2 y_2}=\frac{1}{(y_1-x_1)(x_5-x_1)}=\frac{1}{x_1^2}\cdot\frac{1}{(1-y_1/x_1)(1-x_5/x_1)}.
\]
Expanding in $1/x_1$:
\[
\frac{1}{(1-y_1/x_1)(1-x_5/x_1)}
= 1 + \frac{y_1+x_5}{x_1} + \frac{y_1^2+y_1 x_5+x_5^2}{x_1^2} + O(x_1^{-3}).
\]
So
\[
\frac{1}{x_2 y_2}\bigg|_{1/x_1^4} = \frac{y_1^2+y_1 x_5+x_5^2}{x_1^4}.
\]
The coefficient at $1/x_1^4$ from $a_{x_2 y_2 X_c}/(x_2 y_2 X_c)$ is therefore
\[
\frac{a_{x_2 y_2 X_c}\,(y_1^2+y_1 x_5+x_5^2)}{X_c}.
\]

\emph{The mixed monomial $y_1 x_5/X_c$ is a unique source} for the
coefficient $a_{x_2 y_2 X_c}$, provided $X_c\ne y_1, x_5$ (else degenerate).
That is: $X_c\in F_1\setminus\{y_1,x_5\}=\{x_3,x_4,y_3\}$. So:
\[
\boxed{\text{3 single-term kills at order 4: }
a_{x_2 y_2 x_3}=a_{x_2 y_2 x_4}=a_{x_2 y_2 y_3}=0.}
\]

\paragraph{Higher orders contribute no new single-term kills.}

The monomials with three non-bare/bare-factor combinations either
involve $x_6$ (bare, no tail) or repeat the same Laurent structure already
captured. Direct expansion shows no new isolable single-term equations.

\paragraph{Total kills from $\Z_1$: $35+20+3 = 58$ single-coefficient kills.}

\subsection{Cycling: union over six zones is 151}

By cyclic symmetry every $\Z_r$ produces $35+20+3=58$ kills, with the same
internal structure (just rotated). Tabulating all $6\times 58=348$ raw
kills and accounting for overlaps yields a union of $151$ distinct
coefficients --- exactly the number of non-local multisets at $n=6$.

\paragraph{Sketch of the union argument.} It suffices to show every
non-local multiset is killed by at least one zone. The non-local multisets
at $n=6$ split into types by chord-class composition (\#short, \#long):

\begin{center}
\begin{tabular}{l|c|c}
\toprule
Composition & \# multisets & \# trees among them \\
\midrule
3 short & $\binom{6+2}{3}=56$ & 2 (the two ``zigzag'' triangulations $x_1x_3x_5$, $x_2x_4x_6$) \\
2 short, 1 long & $3\cdot \binom{6+1}{2}=63$ & 6 (six fan-like) \\
1 short, 2 long & $6\cdot \binom{3+1}{2}=36$ & 6 \\
3 long & $\binom{3+2}{3}=10$ & 0 \\
\midrule
Total & 165 & 14 \\
\bottomrule
\end{tabular}
\end{center}

For each non-tree multiset, one identifies a zone $\Z_r$ where the
multiset's indices fit one of the three layers (all in $F_r$;
exactly one non-bare $s$ with the rest in $F_r\setminus\{f_s\}$; or both
non-bare $s$'s with the third in $F_r\setminus\{f_{s_1},f_{s_2}\}$).
A full case analysis confirms every non-tree multiset is caught.

The 14 tree triples are not killed: each is a triangulation (non-crossing
chords), so for every zone the multiset includes either the special $X_*$
itself, or a chord that violates the unique-source conditions of the
three layers.

\begin{theorem}[Locality at six points]
The six hidden one-zeros of the hexagon force every non-local coefficient
of the general rational ansatz $B$ to vanish. The $14$ surviving
coefficients are exactly the tree triangulations of the hexagon.
\end{theorem}

%================================================================
\section{Looking ahead: the structural pattern}

The kill structure at $n=4,5,6$ is governed by three numbers per zone:
\begin{itemize}[nosep,leftmargin=2em]
\item $|F_r|$, the number of free chords (excluding the special).
At $n$ points: $|F_r|=\frac{n(n-3)}{2}-(n-3)-1$.
\item $n-4$, the number of \emph{non-bare} substituted chords.
\item $n-3$, the multiset size (chord factors per monomial of $B$).
\end{itemize}

The general per-zone kill count (verified at $n=4, 5, 6$) is the sum of layers,
where layer $k$ corresponds to multisets with exactly $k$ non-bare substituted
indices:
\[
\text{kills per zone}\;=\;\sum_{k=0}^{\lfloor (n-3)/2\rfloor}
\binom{n-4}{k}\cdot\binom{|F_r|-k+(n-3-k)-1}{n-3-k}.
\]

Plugging in:
\begin{itemize}[nosep,leftmargin=2em]
\item $n=4$: $|F_r|=0$; sum$=0$ (no non-local kills, as expected).
\item $n=5$: $|F_r|=2$; sum$=3+1=4$ kills/zone, union $=10$.
\item $n=6$: $|F_r|=5$; sum$=35+20+3=58$ kills/zone, union $=151$.
\end{itemize}

This is a clean recursive template. The general-$n$ proof should proceed
along the same lines: derive each zero by hand using the mesh identity,
identify the special variable via the enclosed-vertex structure, and verify
layer by layer that the unique-source isolation of single-coefficient kills
covers all non-local multisets. The class hierarchy of chords (short, long,
extra-long, etc.) and the palindromic class structure of the substitutions
suggest an inductive argument that we leave for future work.

\end{document}
